% ===========================================================================
% 2. Background
% ===========================================================================
\section{Background}
\label{sec:background}

\subsection{The DeepMind AI Agent Traps Taxonomy}
\label{sec:bg-taxonomy}

Hammond et al.~\cite{yuntao2025agenttraps} introduce a taxonomy of
environmental attacks targeting AI agents. Unlike adversarial attacks
on model weights or prompts, ``agent traps'' exploit the
\emph{operational environment}---the web pages, APIs, documents, and
inter-agent messages that an autonomous agent must process to complete
its tasks. The taxonomy identifies six attack categories:

\begin{enumerate}[nosep,leftmargin=*]
  \item \category{Content Injection}: Adversarial instructions are
    hidden within otherwise benign content using techniques such as
    invisible CSS styling, HTML comments, image metadata, and dynamic
    cloaking~\cite{greshake2023indirect, yi2023benchmarking}.

  \item \category{Semantic Manipulation}: The framing, emotional
    valence, or social context of content is crafted to bias agent
    decision-making---authority framing, urgency appeals, context
    flooding, and identity manipulation~\cite{perez2022red,
    carlini2024aligned}.

  \item \category{Cognitive State Corruption}: The agent's knowledge
    base or retrieval-augmented generation (RAG) pipeline is poisoned
    through vector embedding manipulation, ranking distortion, gradual
    drift, and cross-contamination~\cite{zou2024poisonedrag,
    zhong2023poisoning, chaudhari2024phantom}.

  \item \category{Behavioural Control}: Deceptive UI elements---
    misleading dialogs, hidden form fields, and infinite loops---trick
    the agent into taking unintended actions~\cite{liao2024eia,
    wu2024agentattack}.

  \item \category{Systemic (Multi-Agent)}: In multi-agent topologies,
    attacks cascade across agents through message poisoning, agent
    impersonation, and cascade failures~\cite{gu2024agent,
    cohen2024here}.

  \item \category{Human-in-the-Loop}: The agent produces biased
    reports that exploit cognitive biases in human
    reviewers---cherry-picked evidence, anchoring effects, and
    decision fatigue~\cite{kahneman2011thinking, tversky1974judgment}.
\end{enumerate}

The original paper provides detailed descriptions and threat models
for each category but stops short of empirical evaluation, noting that
``building reproducible testbeds for these attacks remains an open
challenge''~\cite{yuntao2025agenttraps}.

\subsection{The A2A Protocol}
\label{sec:bg-a2a}

Google's Agent-to-Agent (A2A) protocol~\cite{google2025a2a} defines
a standardized communication layer for autonomous agents. Key
primitives include:

\begin{itemize}[nosep,leftmargin=*]
  \item \textbf{Agent Cards}: JSON-LD discovery documents advertising
    an agent's capabilities, skills, and endpoint URLs.
  \item \textbf{Task lifecycle}: \texttt{message/send} and
    \texttt{message/stream} RPCs that create, update, and complete
    tasks between agents.
  \item \textbf{Artifact exchange}: Typed data payloads (text, files,
    structured data) attached to task messages.
  \item \textbf{Push notifications}: Webhook-based callbacks for
    asynchronous task completion.
\end{itemize}

A2A enables rich multi-agent topologies---hub-and-spoke orchestration,
peer-to-peer delegation, and hierarchical crews---but the protocol
specification provides no built-in authentication, message integrity,
or content validation beyond transport-layer TLS. This makes A2A an
ideal substrate for studying systemic agent traps.

\subsection{Multi-Agent Orchestration}
\label{sec:bg-orchestration}

Modern multi-agent frameworks---AutoGen~\cite{wu2023autogen},
CrewAI~\cite{crewai2024}, LangGraph~\cite{langgraph2024}, and
a2a-crews~\cite{dua2026a2acrews}---enable developers to compose
agents into teams (``crews'') that collaborate on complex tasks.
In these systems, a single compromised agent can propagate corrupted
context to downstream agents, amplifying the blast radius of an
environmental attack. Our testbed specifically targets this
amplification via the A2A protocol.
